\section{Introduction}
\label{sec:intro}

\begin{figure}
\vspace{-0.3cm}
\includegraphics[width=\linewidth]{figures/fig1.pdf}
\vspace{-0.3cm}
\caption{
In this paper, we propose a simple, parameter-free method to uplift any 2D visual features (\emph{e.g.}, CLIP or DINO) into a 3D Gaussian Splatting (3D-GS) 
representation. Uplifting is directly implemented into the rendering process, and takes about 2ms per image and feature dimension. The uplifted features can be leveraged for various 3D tasks, such as segmentation or object removal, based on various inputs, such as scribbles or text prompts.
}
\vspace{-0.2cm}
\label{fig:teaser}
\end{figure}

Image understanding has made remarkable progress, driven by large pretrained models such as CLIP~\citep{radford2021clip}, DINO~\citep{caron2021dino,oquab2024dinov2}, or SAM~\citep{kirillov2023sam}, also called foundation models. A key factor behind their exceptional generalization capabilities lies in the size of their training datasets, often composed of millions or even billions of samples.

Meanwhile, 3D scene representation has also advanced through machine learning approaches like NeRF~\citep{mildenhall2021nerf} and model fitting techniques such as Gaussian Splatting~\citep{kerbl2023gaussiansplatting}. 
These methods typically reconstruct scenes from a few dozen views captured from different angles, effectively preserving both appearance and geometric details. However, they are not suited for semantic tasks.

The complementarity of these two families of approaches has been exploited to integrate geometry and image-level semantics extracted by large 2D pretrained models into NeRF or Gaussian Splatting 3D representations. 
%
This has led to a surge in methods for language-guided object retrieval~\citep{kerr2023lerf, liu2023weakly, zuo2024fmgs, wu2024opengaussian}, scene editing~\citep{kobayashi2022decomposing, chen2024gaussianeditor, fan2023lightgaussian}, or semantic segmentation~\citep{cen2023segmentNeRF, ye2024gaussiangrouping, ying2023omniseg3d} in 3D scenes.

The main limitation of these approaches lies in their dependence on optimization, which requires an iterative process to learn a scene-specific 3D representation by minimizing a reprojection error across all training views. While this loss function is intuitive, a faster and more straightforward method for transferring 2D generic visual features to \emph{already trained} Gaussian Splatting 3D models would be preferable, which is one of the purposes of this work.

In this paper, we demonstrate that a simple, learning-free process is highly effective for uplifting 2D features or semantic masks into 3D Gaussian Splatting scenes. This process, which can be viewed as an `inverse rendering' operation, is both computationally efficient and adaptable to any feature type. We showcase its effectiveness by uplifting visual features from DINOv2~\citep{oquab2024dinov2, darcet2024registers}, semantic masks from SAM~\citep{kirillov2023sam} and SAM2~\citep{ravi2024sam2}, and visual features from CLIP~\citep{ilharco2021openclip}, enabling language-driven applications.
Then, we show that a graph diffusion mechanism~\citep{kondor2002diffusion,smola2003kernels} is helpful for feature refinement in 3D scenes. This mechanism is rooted in spectral graph theory and spectral clustering~\citep{belkin2001laplacian,shi2000normalized,meila2000learning}.
In the context of our work, it transforms coarse segmentation inputs, such as scribbles or alignment scores between visual features and a text query, into accurate 3D segmentation masks without the need for segmentation models such as SAM.
When evaluated on segmentation and open-vocabulary object localization tasks, our method matches state-of-the-art performance while being significantly faster than previous optimization-based approaches.

To summarize, our contributions are threefold: (i)~We introduce a simple, learning-free uplifting approach that can be directly integrated into the rendering process (Sec.~\ref{sec:uplifting}), achieving state-of-the-art segmentation results when applied to SAM-generated semantic masks (Sec.~\ref{sec:method_seg}); 
(ii) We introduce a graph diffusion process for 3D segmentation based on uplifted features (Sec.~\ref{sec:diffusion}) which yields competitive results in both foreground/background segmentation with DINOv2 alone and open-vocabulary segmentation when combined with CLIP (Secs.~\ref{sec:method_seg}, \ref{sec:method_clip});
(iii) We show that our simple uplifting approach can replace more complex and time-consuming feature learning methods and yield strong accuracy gains (Sec.~\ref{sec:method_semantic}).